\chapter{Arbitrage de fusions}\label{ch:s1:merger-arbitrage}

Une entreprise accepte d'être rachetée à 40~\$ par action en numéraire, et son action se négocie à 38.60~\$. Le spread de 1.40~\$, soit 3.6~\%, est ce que le marché facture pour le risque
que l'opération échoue, auquel cas l'action retombe vers les 30~\$ auxquels elle se traitait avant. Acheter à 38.60~\$ revient à vendre une assurance contre cet échec~:
la prime se perçoit par petits montants opération après opération, et les sinistres arrivent ensemble, car les opérations échouent plus souvent quand les marchés baissent. Mitchell et
Pulvino l'ont mesuré sur 4\,750 fusions~: les rendements de l'\omterm{def:s1:merger-arbitrage:arb}{arbitrage de fusions} ne sont pas corrélés au marché lors des phases stables ou haussières, et le sont
lors des fortes baisses, à la manière des rendements de la vente de puts sur indices, et ils dégagent malgré tout des rendements excédentaires d'environ 4~\% par an après coûts, une fois cette non-linéarité contrôlée. Ce chapitre évalue
les spreads des opérations, construit un portefeuille d'opérations synthétiques dont les échecs se concentrent dans les retournements, et mesure ce que rapporte l'assurance. La construction est
\texttt{firm.mergerarb}.

\section{Spreads des opérations et leur évaluation}

\begin{definition}[Arbitrage de fusions, offre publique]\label{def:s1:merger-arbitrage:arb}
L'\emph{arbitrage de fusions}\index{arbitrage de fusions} (merger arbitrage) achète les actions d'une entreprise ayant accepté d'être rachetée, en dessous de la contrepartie offerte, et (dans une
opération en titres) vend à découvert les actions de l'acquéreur qu'il recevra, afin de gagner la différence lorsque l'opération se réalise. Une \emph{offre publique}\index{offre publique} (tender offer) est
une offre adressée directement aux actionnaires d'une entreprise pour acheter leurs actions à un prix déterminé, par opposition à une fusion votée par le conseil d'administration et les actionnaires.
\end{definition}

\begin{definition}[Spread de l'opération, prix en cas d'échec, probabilité implicite de réalisation]\label{def:s1:merger-arbitrage:spread}
Le \emph{spread de l'opération}\index{spread de l'opération} (deal spread) est la contrepartie offerte divisée par le prix de la cible, diminué de un. Le \emph{prix en cas d'échec}\index{prix en cas d'échec} (break price)
est le prix auquel la cible se négocierait si l'opération échouait. La \emph{probabilité implicite de réalisation}\index{probabilité implicite de réalisation} (implied deal probability) est la
probabilité de réalisation qui fait du prix de la cible la moyenne actualisée de la contrepartie et du prix en cas d'échec~: $p = (A - B)/(K - B)$ pour le prix $A$, l'offre $K$
et le prix en cas d'échec $B$, en négligeant l'actualisation.
\end{definition}

Pour l'opération à 40~\$, $p = (38.60 - 30)/(40 - 30) = 0.86$~: le marché intègre une probabilité de 14~\% d'échec. À cette probabilité, la transaction n'a aucun gain espéré
($0.86 \times 1.40 = 0.14 \times 8.60$)~; l'arbitragiste ne gagne que si les opérations échouent moins souvent que ne le suggère le spread. Les spreads sont cotés en rendement annualisé, car
les opérations prennent des mois~: un spread de 3.6~\% sur 84 jours de bourse représente 11.3~\% par an. Les opérations synthétiques (\texttt{firm.mergerarb.simulate\_deals}) sont des opérations en numéraire
annoncées à raison d'une cinquantaine par an sur les dix ans du marché synthétique, avec des primes de 20~\% à 40~\%, et qui se dénouent après un délai médian de 110 jours de bourse~; leur \omterm{def:s1:merger-arbitrage:spread}{prix en cas d'échec}
évolue avec le marché~; leur spread intègre une probabilité d'échec de 12~\%~; leur véritable probabilité d'échec vaut 6~\% plus 1.5 fois la baisse du marché pendant la
durée de l'opération. Le spread moyen à l'entrée est de 4.73~\%, soit 11.8~\% annualisé, et la \omterm{def:s1:merger-arbitrage:spread}{probabilité implicite de réalisation} s'élève à 88~\%.

\begin{dated}{2026-09}{Délais d'attente des fusions aux États-Unis}\label{dat:s1:merger-arbitrage:hsr}
En vertu du programme de notification préalable des fusions Hart-Scott-Rodino, une fois les dossiers des parties complets, celles-ci doivent attendre 30 jours (15 jours pour une \omterm{def:s1:merger-arbitrage:arb}{offre publique} en numéraire
ou une procédure de faillite) avant de réaliser l'opération, sauf si les autorités accordent une résiliation anticipée~; après un examen préliminaire, l'autorité chargée de l'examen (la FTC ou le
département de la Justice) peut laisser le délai expirer, y mettre fin par anticipation, ou émettre une demande de renseignements complémentaires (Second Request) qui prolonge l'examen.
\end{dated}

\section{Risque d'échec et corrélation avec le marché}

\begin{center}
\small
\setlength{\tabcolsep}{4pt}%
\begin{tabular}{@{}p{5.4cm}cccc@{}}
\toprule
portefeuille d'opérations synthétiques, années 3 à 10 & cas de base & échecs & spread intégré & les deux\\
 & & indépendants & à 6~\% & \\
\midrule
rendement annuel~; au-dessus du taux de 4~\% & 6.76~\%~; 2.76~\% & ---~; 4.30~\% & ---~; $-0.47~\%$ & ---~; 1.03~\%\\
volatilité~; ratio de Sharpe (excédentaire) & 3.26~\%~; 0.85 & ---~; 1.50 & ---~; $-0.17$ & ---~; 0.46\\
bêta dans les mois où le marché a baissé de plus de 4~\% & 0.20 & 0.20 & 0.17 & 0.17\\
bêta dans les autres mois & 0.09 & 0.11 & 0.04 & 0.06\\
opérations menées à terme & 92.5~\% & 94.7~\% & 92.5~\% & 94.7~\%\\
\bottomrule
\end{tabular}
\end{center}

Le portefeuille de référence détient 22 opérations en moyenne, 374 sur les huit années, pondérées de manière égale. Il rapporte 2.76~\% par an au-dessus du taux auquel ses spreads sont actualisés,
avec une volatilité de 3.26~\% et un ratio de Sharpe de 0.85. Les opérations ont abouti dans 92.5~\% des cas, contre les 88~\% qu'impliquaient les spreads~: l'écart est la prime. Les opérations
ont échoué dans 13.7~\% des cas lorsque le marché a baissé pendant leur durée (117 opérations), et dans 4.7~\% des cas lorsqu'il a monté. Le bêta du portefeuille vaut 0.20 durant les dix mois où le
marché a baissé de plus de 4~\%, et 0.09 durant les 86 autres~: un profil de type put, plus faible que dans les données historiques parce que le marché synthétique avait peu de
fortes baisses. Deux colonnes contrefactuelles séparent les composantes. Si les échecs ne dépendaient pas du marché, les mêmes spreads rapporteraient 4.30~\%~: la concentration des échecs
dans les retournements coûte 1.5 point par an. Si les spreads n'intégraient que le taux d'échec en marché stable, le portefeuille perdrait 0.47~\% par an~: la prime
contenue dans le spread est précisément le prix du risque de krach.

\begin{omfigure}
\begin{tikzpicture}
\begin{axis}[width=8.5cm,height=6.0cm,xlabel={\small rendement du marché dans le mois (\%)},ylabel={\small rendement du portefeuille d'opérations (\%)},
  tick label style={font=\footnotesize},xmin=-16,xmax=12,ymin=-4,ymax=3.5,grid=major,grid style={gray!25},table/col sep=comma]
\addplot[omDef,only marks,mark=*,mark size=1.4pt] table[x=market_pct,y=deals_pct]{figdata/strategies-1/11-merger-arbitrage/months.csv};
\addplot[black!40,dashed] coordinates {(-4,-4) (-4,3)};
\end{axis}
\end{tikzpicture}

\omcaption{Rendements mensuels du portefeuille d'opérations synthétiques face à ceux du marché, années 3 à 10~; la ligne en pointillés marque une baisse du marché de 4~\%. Données~:
\texttt{s1\_mergerarb.run}.}\label{fig:s1:merger-arbitrage:months}
\end{omfigure}

Le pire mois du portefeuille ($-3.06~\%$) est un mois au cours duquel le marché a monté de 2.8~\%~: avec vingt opérations, deux ou trois échecs simultanés suffisent, et
les échecs propres à chaque opération représentent l'autre moitié du risque. Un livre d'\omterm{def:s1:merger-arbitrage:arb}{arbitrage de fusions} diversifie cette moitié en détenant de nombreuses opérations~; il ne peut pas diversifier la moitié qui
provient du marché.

\section{Opérations en titres et collars}

\begin{definition}[Collar d'une offre en titres]\label{def:s1:merger-arbitrage:collar}
Un \emph{collar d'une offre en titres}\index{collar d'une offre en titres} (stock-deal collar) fixe le rapport d'échange d'une opération en titres tant que le cours de l'acquéreur reste à l'intérieur d'une fourchette, et fixe la
valeur de la contrepartie en dehors de celle-ci (le rapport s'ajuste alors), protégeant les actionnaires de la cible contre les fortes baisses du cours de l'acquéreur.
\end{definition}

Dans une opération en titres, les détenteurs de la cible reçoivent un nombre d'actions de l'acquéreur~; l'arbitragiste achète la cible et vend à découvert ce nombre d'actions de l'acquéreur,
de sorte que le spread ne dépend plus du cours de l'acquéreur. Un collar modifie la couverture. Avec un rapport de 0.5 et une fourchette de 40~\$ à 50~\$, un cours de l'acquéreur
à 42~\$ donne une contrepartie de 21~\$, couverte par la vente à découvert de 0.5 action de l'acquéreur~; en dessous de 40~\$, à 36~\$, la contrepartie est fixée à 20~\$ et la couverture est nulle, puisque
la valeur ne bouge plus avec l'acquéreur. Le ratio de couverture est le delta du collar (\texttt{collar\_hedge}), et près des bornes de la fourchette il fait un saut~: les opérations sous collar
doivent être rééquilibrées, et leurs spreads portent le gamma de cette option.

\section{Risque réglementaire et risque de financement}

Les opérations échouent pour des raisons que le niveau du marché ne capture pas. Le contrôle antitrust (les délais d'attente de l'encadré daté, et une Second Request qui peut les prolonger
de plusieurs mois) est la mieux connue~; un financement qui disparaît, un vote d'actionnaires qui échoue, une offre concurrente (qui fait au contraire monter le prix) et des changements défavorables importants
dans l'activité de la cible sont les autres. Chacun a sa propre probabilité et son propre \omterm{def:s1:merger-arbitrage:spread}{prix en cas d'échec}, et le spread d'une opération est un mélange de ces éléments. Le
modèle synthétique les regroupe en une seule probabilité, dont une part dépend du marché~; un desk réel les estime opération par opération, lit l'accord de fusion, et dimensionne les positions
selon ce que chaque opération peut faire perdre, et non selon son spread.

\section{Portefeuille d'opérations}

L'\omterm{def:s1:merger-arbitrage:arb}{arbitrage de fusions} est limité par le capital autant que par les opportunités. Baker et Savasoglu ont étudié les limites de ce capital dans
les fusions et acquisitions~; la rémunération de type put de Mitchell et Pulvino en est une des raisons. Le budget de risque d'un livre d'opérations comporte trois postes~: les échecs propres à chaque opération,
qui se diversifient sur vingt opérations ou plus~; les échecs dépendant du marché, qui ne se diversifient pas~; et l'exposition au marché des \omterm{def:s1:merger-arbitrage:spread}{prix en cas d'échec} eux-mêmes, qui
apparaît comme le bêta du livre lors des baisses. Les deux derniers constituent l'assurance vendue. Le véritable indice de référence du livre est une vente de put à découvert sur l'indice, et non le cash.

\section{Fiches de stratégie}

\begin{strategyfile}{Spread d'une opération en espèces}\label{strat:s1:merger-arbitrage:cash}
\sfield{Qui vous paie, et pourquoi} Les actionnaires de la cible qui vendent après l'annonce plutôt que d'attendre et de supporter le risque d'échec.
\sfield{Instruments et plateformes} Les actions de la cible.
\sfield{Signal} Le spread annualisé par rapport à une estimation de la probabilité d'échec et du \omterm{def:s1:merger-arbitrage:spread}{prix en cas d'échec}.
\sfield{Dimensionnement et exécution} Dimensionner selon la perte en cas d'échec de l'opération, et non selon le spread~; entrer après l'annonce.
\sfield{Coûts} Faible rotation~; financement de la position.
\sfield{Comment elle s'éteint} Les échecs, surtout en période de repli~; les surprises réglementaires.
\sfield{Horizon, capacité, infrastructure} Quelques mois~; capacité limitée par la taille de l'opération~; analyse juridique et réglementaire.
\sfield{Backtest honnête} Opérations telles qu'annoncées, y compris celles qui ont échoué~; \omterm{def:s1:merger-arbitrage:spread}{prix en cas d'échec} estimés avant l'issue.
\sfield{Sources} Mitchell et Pulvino (2001)~: 4\,750 fusions, 1963--1998, rendements en excès d'environ 4~\% par an après coûts, rémunération de type put.
\end{strategyfile}

\begin{strategyfile}{Spread stock-for-stock couvert}\label{strat:s1:merger-arbitrage:stock}
\sfield{Qui vous paie, et pourquoi} Comme pour les opérations en espèces, avec le risque de prix de l'acquéreur couvert.
\sfield{Instruments et plateformes} La cible, en position longue~; l'acquéreur, en position courte au ratio d'échange.
\sfield{Signal} L'écart entre le cours de la cible et le ratio multiplié par le cours de l'acquéreur.
\sfield{Dimensionnement et exécution} Maintenir la position courte au ratio d'échange~; ajuster pour les dividendes.
\sfield{Coûts} Frais d'emprunt sur l'acquéreur, parfois élevés lorsque de nombreux arbitragistes sont vendeurs.
\sfield{Comment elle s'éteint} Les échecs~; un short squeeze sur l'acquéreur si beaucoup sont vendeurs.
\sfield{Horizon, capacité, infrastructure} Quelques mois~; disponibilité des titres à emprunter.
\sfield{Backtest honnête} Frais d'emprunt et rappels sur l'acquéreur~; mouvement de l'acquéreur en cas d'échec.
\sfield{Sources} Mitchell et Pulvino (2001).
\end{strategyfile}

\begin{strategyfile}{Opérations à collar}\label{strat:s1:merger-arbitrage:collar}
\sfield{Qui vous paie, et pourquoi} Comme pour les opérations en titres, avec une option incorporée dans la contrepartie.
\sfield{Instruments et plateformes} La cible, l'acquéreur, et éventuellement des options cotées pour couvrir le gamma du collar.
\sfield{Signal} L'écart par rapport à la valeur du collar.
\sfield{Dimensionnement et exécution} Un ratio de couverture égal au delta du collar, rééquilibré près des bornes de la bande.
\sfield{Coûts} Rééquilibrage de la couverture.
\sfield{Comment elle s'éteint} Erreurs de couverture près des bornes de la bande~; échecs.
\sfield{Horizon, capacité, infrastructure} Quelques mois~; un modèle d'options pour le collar.
\sfield{Backtest honnête} Les conditions du collar telles qu'écrites, y compris les périodes de moyenne.
\sfield{Sources} Le calcul du collar de ce chapitre~; aucun historique de performance vérifié.
\end{strategyfile}

\begin{strategyfile}{Portefeuille diversifié d'opérations}\label{strat:s1:merger-arbitrage:portfolio}
\sfield{Qui vous paie, et pourquoi} Le marché, pour la prise en charge d'un risque d'échec concentré lors des replis.
\sfield{Instruments et plateformes} Vingt opérations annoncées ou plus.
\sfield{Signal} Le spread de chaque opération par rapport à son risque estimé.
\sfield{Dimensionnement et exécution} Risque égal par opération~; limites sur les opérations d'un même secteur ou soumises au même régulateur.
\sfield{Coûts} Faibles.
\sfield{Comment elle s'éteint} Les krachs boursiers, lorsque les échecs se regroupent et que les \omterm{def:s1:merger-arbitrage:spread}{prix en cas d'échec} baissent ensemble.
\sfield{Horizon, capacité, infrastructure} Continu~; capacité fixée par le flux des opérations.
\sfield{Backtest honnête} Toutes les opérations, en données disponibles à chaque date~; comparaison avec la vente d'un put à découvert sur l'indice.
\sfield{Sources} Mitchell et Pulvino (2001)~; simulation de ce chapitre (2.76~\% par an au-dessus du taux, bêta de 0.20 les mois de baisse et de 0.09 sinon).
\end{strategyfile}

\section{Tutoriel~: vendre de l'assurance}\label{tut:s1:merger-arbitrage:tut}

\begin{tutorial}
\textbf{Objectif.} Tarifer les spreads d'opérations, simuler un portefeuille d'opérations dont les échecs dépendent du marché, et mesurer son gain. \textbf{État final~:} le tableau
et la \cref{fig:s1:merger-arbitrage:months}.
\begin{tutsteps}
\item \textbf{Tarification}~: probabilité implicite et collar.
\omcode{code/firm/mergerarb/firm_mergerarb.py}{36}{48}{Probabilité implicite de réalisation et conditions du collar.}\label{lst:s1:merger-arbitrage:pricing}
\item \textbf{Opérations}~: annonce, trajectoire des cours, échec ou réalisation.
\omcode{code/firm/mergerarb/firm_mergerarb.py}{65}{88}{Vie d'une opération, avec échecs dépendant du marché.}\label{lst:s1:merger-arbitrage:deals}
\item \textbf{Exécuter} \texttt{run()} et ses contrefactuels, \texttt{break\_rates()} et \texttt{fig\_mergerarb.py}.
\end{tutsteps}
\textbf{Pour aller plus loin.} Ajouter des opérations en titres assorties de collars, et les couvrir~; laisser les spreads s'élargir après les baisses du marché (le marché réévalue $q$)~; plafonner
l'exposition du livre aux opérations d'un même secteur.
\end{tutorial}

\section{Construction~: arbitrage de fusions}\label{bld:s1:merger-arbitrage:mergerarb}

\begin{build}
\bfield{Rôle} Tarification des spreads d'opérations et portefeuille d'opérations à échecs dépendant du marché.
\bfield{Interface} \texttt{spread(price, offer)}, \texttt{annualised(spread, days)}, \texttt{implied\_probability(price, offer, break\_price, rate, days)},
\texttt{collar\_value(acquirer, ratio, low, high)}, \texttt{collar\_hedge(acquirer, ratio, low, high)}, \texttt{DealConfig}, \texttt{simulate\_deals(mkt, cfg, rng)},
\texttt{portfolio(deals, T)}.
\bfield{Règles} Opérations entrées après leur annonce~; échecs décidés par l'information disponible à la fin de l'opération~; rendements en excès du taux d'actualisation rapportés.
\bfield{Tests d'acceptation} \texttt{code/firm/mergerarb/tests/}~: spread, probabilité implicite et annualisation calculés à la main~; valeur et couverture du collar~; opérations sur un marché plat
qui échouent au taux de base, et le portefeuille gagne un rendement positif.
\bfield{Pour aller plus loin} Opérations en titres assorties de couvertures par l'acquéreur~; offres concurrentes~; spreads qui s'élargissent lorsque le marché baisse.
\end{build}

\begin{omsources}
\item M.~Mitchell et T.~Pulvino, ``Characteristics of risk and return in risk arbitrage'', \emph{Journal of Finance} 56(6), 2001.
\item M.~Baker et S.~Savasoglu, ``Limited arbitrage in mergers and acquisitions'', \emph{Journal of Financial Economics} 64(1), 2002.
\item US Federal Trade Commission, Premerger Notification and the Merger Review Process.
\end{omsources}

\section{Exercices}

\begin{exercise}[$\star$]\label{exo:s1:merger-arbitrage:1}
Une offre en espèces de 40~\$ se négocie à 38.60~\$, avec un \omterm{def:s1:merger-arbitrage:spread}{prix en cas d'échec} de 30~\$. Calculer le spread et la \omterm{def:s1:merger-arbitrage:spread}{probabilité implicite de réalisation}.
\end{exercise}

\begin{exercise}[$\star$]\label{exo:s1:merger-arbitrage:2}
L'opération doit se clore dans 84 jours de bourse. Calculer l'annualisation du spread.
\end{exercise}

\begin{exercise}[$\star$]\label{exo:s1:merger-arbitrage:3}
Une opération en titres verse 0.5 action de l'acquéreur, avec un collar entre 40~\$ et 50~\$. Calculer la contrepartie pour des cours de l'acquéreur de 42~\$ et de 36~\$, puis le
nombre d'actions de l'acquéreur à vendre à découvert par action de la cible dans chaque cas.
\end{exercise}

\begin{exercise}[$\star\star$]\label{exo:s1:merger-arbitrage:4}
Montrer que l'achat, à 38.60~\$, d'une opération à 40~\$ a un profit espéré nul à la probabilité implicite, et préciser ce que l'arbitragiste doit croire pour trader.
\end{exercise}

\begin{exercise}[$\star\star$]\label{exo:s1:merger-arbitrage:5}
Expliquer pourquoi le bêta d'un portefeuille d'opérations augmente sur les marchés en baisse, même si ses opérations n'ont aucun lien entre elles.
\end{exercise}

\begin{exercise}[$\star\star$]\label{exo:s1:merger-arbitrage:6}
Expliquer ce que fait une Second Request au spread d'une opération, et pourquoi.
\end{exercise}

\begin{exercise}[$\star\star\star$]\label{exo:s1:merger-arbitrage:7}
\emph{Programmation.} Exécuter \texttt{run(0.0, 0.12, 0.06)} et \texttt{run(1.5, 0.06, 0.06)}. Expliquer ce que isole chaque contrefactuel.
\end{exercise}

\begin{exercise}[$\star\star\star$]\label{exo:s1:merger-arbitrage:8}
\emph{Trouver la faille.} «~L'\omterm{def:s1:merger-arbitrage:arb}{arbitrage de fusions} affiche un ratio de Sharpe proche de un et un bêta proche de zéro~; il est neutre au marché.~»
\end{exercise}

\section{Problème~: vendre de l'assurance}

\begin{problem}[{Problème du week-end --- la prime et les sinistres}]\label{pb:s1:merger-arbitrage:1}
Les opérations synthétiques du chapitre et les données publiques.

\medskip\noindent\textbf{Partie I --- Tarification.}
\begin{enumerate}
\item Définir l'\omterm{def:s1:merger-arbitrage:arb}{arbitrage de fusions}, une \omterm{def:s1:merger-arbitrage:arb}{offre publique}, le \omterm{def:s1:merger-arbitrage:spread}{spread de l'opération} et le \omterm{def:s1:merger-arbitrage:spread}{prix en cas d'échec}.
\item Calculer la \omterm{def:s1:merger-arbitrage:spread}{probabilité implicite de réalisation} de l'opération initiale.
\item Pourquoi les spreads sont-ils annualisés~?
\item Décrire les opérations synthétiques.
\end{enumerate}

\medskip\noindent\textbf{Partie II --- Le portefeuille.}
\begin{enumerate}[resume]
\item Donner son rendement, son rendement en excès, sa volatilité et son ratio de Sharpe.
\item Comparer les taux de réalisation implicite et réalisé.
\item Comment les taux d'échec dépendent-ils du marché~?
\item Donner les bêtas dans les mois de baisse et dans les autres mois.
\end{enumerate}

\medskip\noindent\textbf{Partie III --- Les composantes.}
\begin{enumerate}[resume]
\item Que montrent les deux contrefactuels~?
\item Décrire le pire mois et ce qu'il révèle sur le risque spécifique à l'opération.
\item Comment couvre-t-on les opérations en titres et les collars~?
\item Quelles sont les raisons non liées au marché qui font échouer les opérations~?
\end{enumerate}

\medskip\noindent\textbf{Partie IV --- Le verdict.}
\begin{enumerate}[resume]
\item Énoncer le \emph{résultat nommé}~: le bêta du portefeuille dans les marchés en baisse et en hausse, et sa \omterm{def:s1:merger-arbitrage:spread}{probabilité implicite de réalisation} face à la réalisation effective.
\item Que trouvent Mitchell et Pulvino~?
\item Quel est le bon indice de référence pour un portefeuille d'opérations~?
\item Comment dimensionneriez-vous les opérations individuelles~?
\item Que signifie le délai d'attente HSR pour le calendrier d'une opération~?
\item Quel risque un portefeuille d'opérations peut-il diversifier, et lequel ne le peut-il pas~?
\item Quand ajouteriez-vous au portefeuille~: après une baisse du marché ou après une hausse~?
\item En une phrase~: que vend un arbitragiste de fusions~?
\end{enumerate}
\end{problem}

\section{Questions d'entretien}

\begin{interviewq}[$\star$ \iqroles{trader, chercheur}]\label{iq:s1:merger-arbitrage:1}
Qu'est-ce que l'\omterm{def:s1:merger-arbitrage:arb}{arbitrage de fusions}, et d'où provient son rendement~?
\end{interviewq}

\begin{interviewq}[$\star\star$ \iqroles{trader}]\label{iq:s1:merger-arbitrage:2}
Comment couvririez-vous une opération en titres contre titres avec un collar~?
\end{interviewq}

\begin{interviewq}[$\star\star$ \iqroles{risque}]\label{iq:s1:merger-arbitrage:3}
Pourquoi compare-t-on un portefeuille d'\omterm{def:s1:merger-arbitrage:arb}{arbitrage de fusions} à un put vendu~?
\end{interviewq}

\begin{interviewq}[$\star\star$ \iqroles{chercheur}]\label{iq:s1:merger-arbitrage:4}
Comment estimeriez-vous la probabilité d'échec d'une opération et son \omterm{def:s1:merger-arbitrage:spread}{prix en cas d'échec}~?
\end{interviewq}

\begin{interviewq}[$\star\star$ \iqroles{trader}]\label{iq:s1:merger-arbitrage:5}
Le spread d'une opération double du jour au lendemain sans aucune nouvelle. Que peut-il s'être passé~?
\end{interviewq}

\begin{interviewq}[$\star\star\star$ \iqroles{chercheur}]\label{iq:s1:merger-arbitrage:6}
Établir la \omterm{def:s1:merger-arbitrage:spread}{probabilité implicite de réalisation} en actualisant au taux $r$ sur $\tau$ années, et montrer comment le rendement en excès espéré de l'arbitragiste dépend
de l'écart entre la probabilité implicite et la probabilité réelle.
\end{interviewq}
